\uuid{ve5X}
\titre{Polynôme/Formule de Taylor}
\auteur{hanani ; mzari}
\organisation{exo7}
\licence{CC BY-NC-SA 4.0}
\datecreate{2020-01-20}
\niveau{L1}
\module{Algèbre}
\chapitre{Déterminant, système linéaire}
\sousChapitre{Système linéaire, rang}
\difficulte{3}
\theme{Systèmes d'équations linéaires, Polynôme, Formule de Taylor}
\exo7qcmid{lille-2.systemes.026}

\contenu{
\begin{question}
\qclassification{200.03}
\qtags{motcle=Polynôme/Formule de Taylor}
Soit $P$ un polynôme à coefficients réels de degré $\le 11$ vérifiant les conditions :
$$P(1)=1!,\quad P'(1)=2!,\quad P''(1)=3!, \dots , \quad
P^{(10)}(1)= 11!.$$
Quelles sont les assertions vraies ?
\begin{answers}
\bad{Un tel polynôme $P$ n'existe pas}.
\good{Il existe une infnité de polynômes $P$ vérifiant ces conditions}.
\bad{Il existe un unique polynôme $P$ vérifiant ces conditions}.
\bad{Si $P$ est un polynôme qui vérifie ces conditions, alors
$P(X)=1+2(X-1)+ 3(X-1)^2+ \dots + 11(X-1)^{10} +  12(X-1)^{11}$}.
\end{answers}
\begin{explanations} On peut résoudre un système dont les inconnues ($12$ inconnues) sont les coefficients du polynôme, mais cela est long ! Par contre, en appliquant la formule de Taylor à $P$, au voisinage de $1$,  on a :
$$P(X)=P(1)+(X-1)\frac{P'(1)}{1!} + (X-1)^2\frac{P''(1)}{2!}+ \dots + (X-1)^{11}\frac{P^{(11)}(1)}{11!}.$$
En utilisant les conditions que doit vérifier $P$, on obtient :
$$P(X)=1+2(X-1)+ 3(X-1)^2+ \dots + 11(X-1)^{10} +  a(X-1)^{11},$$
où $a$ est un réel. Par conséquent, il existe une infinité de polynômes vérifiant les conditions ci-dessus.
\end{explanations}
\end{question}
}
